% Sketches of schwartz-space-n-fold-cover-sl2.tex, keyed by its labels; see ~/work/sketches.
% Most sentences are taken or condensed from the paper.  Drafts awaiting the author's check.

\begin{sketch}[document]
  Poisson summation on the plane gives the functional equation for Eisenstein series on $\mathrm{SL}_2$.  The paper develops the analogue for Kubota's $n$-fold cover, in the spirit of Braverman--Kazhdan.  Locally, the Schwartz space consists of genuine functions on Deligne's $\mu_n$-torsor over the punctured plane, and the Fourier transform $\mathcal{F}$ is defined like the symplectic Fourier transform, with the symplectic form replaced by its $n$th roots.  At almost all places a basic vector is fixed by $\mathcal{F}$.  Globally,
  \[
    \sum_{v \in F^2 - \{0\}} \Phi(v) + B(b \Phi)
    =
    \sum_{v \in F^2 - \{0\}} \mathcal{F} \Phi(v) + B(b \mathcal{F} \Phi),
  \]
  where the leading terms $b \Phi$ and $b \mathcal{F} \Phi$ at the origin take the place of the values $\Phi(0)$ and $\mathcal{F} \Phi(0)$.  Taking a Fourier coefficient gives a Voronoi-type summation formula for Gauss sums of order $n$, with an explicit integral transform relating the weights.  For cubic Gauss sums, it takes sextic residue twists to quadratic twists, as in recent work on the quasi-Riemann hypothesis.
\end{sketch}

% Section 1

\begin{sketch}[sec:idea-of-the-proof]
  For $g \in \widetilde{\mathrm{SL}}_2(\mathbb{A})$, the difference of the sums of $g \Phi$ and of $g \mathcal{F} \Phi$ over $F^2 - \{0\}$ is an automorphic function of $g$.  Ordinary Poisson summation on $\mathbb{A}$ and on $\mathbb{A}^\times / \mathbb{A}^{\times n}$ gives its constant term in terms of $b \Phi$ and $b \mathcal{F} \Phi$.  If each of these leading terms has a zero factor, the constant term vanishes, the difference decays rapidly in the cusp, and unfolding shows that it is orthogonal to itself, hence zero.  For general $\Phi$, dilation at one finite place scales the two leading terms differently, which separates their contributions.
\end{sketch}

% Sections 3-4

\begin{sketch}[sec:genuine]
  Deligne's torsor is glued from the trivial torsors on $U_X$ ($x \neq 0$) and $U_Y$ ($y \neq 0$) by the Hilbert symbol.  A genuine function is a pair $(\Phi_X, \Phi_Y)$ of smooth functions on the two charts with
  \[
    \Phi_X(x, y) = \{x, y\}_n \Phi_Y(x, y)
  \]
  where both are defined.  The test functions $\mathcal{D}$ are small near the origin and near infinity: compactly supported in $U$ at a finite place, and rapidly decaying at both ends, with all derivatives, at an archimedean place.  For $n = 1$, a genuine function is a function on the punctured plane.
\end{sketch}

\begin{sketch}[definition:lines]
  A genuine function is not a function, so integrating it over a line requires a trivialization of the torsor along the line.  Deligne uses his trivialization along $a x + b y = 1$ to realize the standard intertwining operator as a Radon transform; dilation by $\tau^n$ transports it to $a x + b y = \tau^n$.  With $(a, b) = (y', -x')$, this gives the integrals $R_X \Phi(v', \tau)$ and $R_Y \Phi(v', \tau)$ over the line $x y' - y x' = \tau^n$.  The Hilbert symbol factors make the integrands continuous where the line crosses the axes, and the two integrals differ by the factor $\{x', y'\}_n$ (Lemma \ref{lemma:lines}).
\end{sketch}

\begin{sketch}[definition:fourier-transform]
  Integrate over the line $x y' - y x' = \tau^n$, then take the Fourier transform in $\tau$ at $1$; the factor $\lvert n \rvert^{-1/2}$ makes the result unitary.  For $n = 1$ this is the symplectic Fourier transform.  Equivalently (Lemma \ref{lemma:root-kernel}), $\mathcal{F}$ has the kernel
  \[
    K(w, v) = \{x, -x'\}_n \, \lvert n \rvert^{-3/2} \sum_{\tau^n = [v, w]} \psi(\tau) \, \lvert \tau \rvert^{1-n},
  \]
  the sum running over the $n$th roots of the symplectic form $[v, w] = x y' - y x'$.

  \emph{Why the transform in $\tau$?}  On a function $f$ in the principal series, the line integrals satisfy $R f(v', \tau) = \xi(\tau)^{-1} \lvert \tau \rvert^{-s} A f(v')$, where $A$ is the standard intertwining operator.  The transform in $\tau$ therefore multiplies $A f$ by Tate's factor $\gamma(s, \xi, \psi)$.  At an unramified place this supplies the factor $L(1 - s) / L(s)$ that $A$ lacks, so that $\mathcal{F}$ carries the Mellin coefficient $\Phi^0[s] = L(1 + s) e_s$ of the basic vector to $\Phi^0[-s]$, as the Fourier transform of the plane does for the characteristic function of $\mathfrak{o}^2$.
\end{sketch}

\begin{sketch}[definition:schwartz-space]
  The space $\mathcal{D}$ is not preserved by $\mathcal{F}$, so, following Braverman and Kazhdan, one enlarges it by the transforms of its elements.  The result is preserved by the dilations by $F^{\times n}$ and does not depend on $\psi$.  At a finite place, an element of $\mathcal{S}$ vanishes outside a compact subset of $F^2$ and agrees near the origin with its leading term, which is $A f$ for a homogeneous genuine function $f$ (Corollary \ref{corollary:finite-schwartz}).  For $n = 1$, $\mathcal{S}$ consists of the restrictions to $U$ of the Schwartz--Bruhat functions on $F^2$.
\end{sketch}

% Section 6

\begin{sketch}[proposition:basic-vector-fixed]
  At a finite place prime to $n$, the basic vector replaces the characteristic function of $\mathfrak{o}^2$.  It is supported on the vectors $\varpi^k (x_0, y_0)$ with $(x_0, y_0)$ primitive and $k \in n \mathbb{Z}_{\geq 0}$, where it has absolute value $q^{(n-1) k / n}$.  This weight makes its Mellin coefficient $L(1 + s) e_s$, as for $n = 1$, and the values at the primitive vectors are chosen so that it is invariant under the unipotent translations by $\mathfrak{o}$.  The proof of $\mathcal{F} \Phi^0 = \Phi^0$ computes $\mathcal{F} \Phi^0$ on the line $x = 0$ and uses that invariance.  The primitive layer $\Phi^1 = \Phi^0 - q^{n-1} \Phi^0(\varpi^{-n} \, \cdot)$ lies in $\mathcal{D}$, and the identity
  \[
    (1 - q^{-2}) \, \Phi^0 = \mathcal{F} \Phi^1 + q^{-n-1} \, \Phi^1(\varpi^n \, \cdot)
  \]
  puts $\Phi^0$ in $\mathcal{S}$.
\end{sketch}

% Section 7

\begin{sketch}[theorem:schwartz-space]
  For $n = 1$, the symplectic Fourier transform is an involution commuting with $\mathrm{SL}_2(F)$.  Here $\mathcal{F}$ is a unitary bijection of $\mathcal{S}$ commuting with the metaplectic group, and $\mathcal{F}^2 = h(\varepsilon)$ with $\varepsilon = (-1)^{n+1}$; for odd $n$, $h(\varepsilon)$ is the identity.  The sign $\varepsilon$ has two sources: Fourier inversion in $\tau$ replaces $\tau^n$ by $(-1)^n \tau^n$, and the antisymmetry of $x y' - y x'$ contributes another.

  The inversion formula is proved in Appendix \ref{sec:marginals}.  There $\mathcal{F}$ exchanges the integrals of $\Phi_X$ along the lines $x = \mathrm{const}$ with the values of $\Phi_Y$ on the line $x = 0$, each time by a Fourier transform in an $n$th root variable, so $\mathcal{F}^2$ acts on the axis values by Fourier inversion.  The translations $n^-(u)$ commute with $\mathcal{F}$ and move every point of $U_Y$ onto that axis.  Unitarity follows from a symmetry of the kernel $K$.
\end{sketch}

% Section 8

\begin{sketch}[sec:germs]
  For $n \geq 2$, an element $\Psi$ of $\mathcal{S}$ has no natural value at the origin.  Its role is played by the leading term
  \[
    b \Psi(u) = \lim_{\mu \to 0} \lvert \mu \rvert^{n-1} \Psi(\mu^n u),
  \]
  a genuine function on $U$ with $(b \Psi)(\mu^n u) = \lvert \mu \rvert^{1-n} \, b \Psi(u)$.  It vanishes on $\mathcal{D}$, and at a finite place the kernel of $b$ on $\mathcal{S}$ is exactly $\mathcal{D}$.  For $\Phi \in \mathcal{D}$, $b \mathcal{F} \Phi$ is $\lvert n \rvert^{-1/2}$ times the integral over $\tau$ of the line integrals; for $n = 1$ this is $\int_{F^2} \Phi = \mathcal{F} \Phi(0)$.  It is a multiple of $A(P \Phi)$, where $P \Phi$ averages $\Phi$ over the dilations by $F^{\times n}$, so the vanishing condition $P \Phi = 0$ forces $b \mathcal{F} \Phi = 0$.
\end{sketch}

% Section 9

\begin{sketch}[theorem:poisson]
  The global Schwartz space is the restricted tensor product of the local ones with respect to the basic vectors.  By Hilbert's reciprocity law, the local values of a pure tensor multiply to well-defined values $\Phi(v)$ at the points of $F^2 - \{0\}$.  The boundary functional $B$ plays the role of $\Phi \mapsto \Phi(0)$: for $n = 1$ it is $\prod_{\mathfrak{p}} \Phi_{\mathfrak{p}}(0)$ on pure tensors, and the formula is the usual Poisson summation formula.  For $n \geq 2$, $B$ exists and is determined by the formula, but is not computed.  The proof is in Appendix \ref{sec:proof-poisson}.

  \emph{When does $B$ drop out?}  When some factor of $\Phi$ and some factor of $\mathcal{F} \Phi$ have leading term zero (Corollary \ref{corollary:poisson-vanishing}), for instance when some $\Phi_{\mathfrak{p}} \in \mathcal{D}$ satisfies the vanishing condition.
\end{sketch}

% Section 10

\begin{sketch}[theorem:main]
  Take the Fourier coefficient of the sum $\sum_v \Phi(v)$ along the translations $(x, y) \mapsto (x, y + u x)$, against $\psi(-u)$.  The terms with $x = 0$ drop out, and those with $x = c$ unfold to products of the local Jacquet integrals
  \[
    W_{\mathfrak{p}}(c) = \int_{F_{\mathfrak{p}}} \Phi_{\mathfrak{p}, X}(c, y) \, \psi_{\mathfrak{p}}(-y / c) \, d y.
  \]
  Every test function on $F_{\mathfrak{p}}^\times$ arises in this way, and $\mathcal{F}$ induces a transform $W_{\mathfrak{p}} \mapsto W_{\mathfrak{p}}^\vee$ (Proposition \ref{proposition:transform-formula}).  Outside $S$, the basic vector gives the basic function, a generating function for Gauss sums of order $n$, fixed by the transform when $\psi_{\mathfrak{p}}$ is unramified.  Applying this to the two sides of the Poisson formula gives
  \[
    \sum_{c \in F^\times} W(c) = \sum_{c \in F^\times} W^\vee(c) + \textrm{main term}.
  \]
  The main term is the Fourier coefficient of the boundary term.  It is a linear combination of the products over $\mathfrak{p} \in S$ of the integrals of $W_{\mathfrak{p}}(x) \lvert x \rvert_{\mathfrak{p}}^{1/n}$ over the cosets of $F_{\mathfrak{p}}^{\times n}$, with coefficients $\rho_S(\beta)$ that are not determined.  It vanishes for $n = 1$, and whenever these integrals all vanish at one place.
\end{sketch}

\begin{sketch}[proposition:transform-formula]
  As for $\mathcal{F}$, the transform is a Fourier transform in an $n$th root variable: an integral of $W$ against a Hilbert symbol and $\psi(-t^n / c x)$, followed by the Fourier transform in $t$ at $1$.  For $n = 1$ it is $W^\vee(c) = W(1/c)$.  This explicit integral is the definition of $W^\vee$ used in the Lean statement of Theorem \ref{theorem:gauss-sums}.
\end{sketch}

\begin{sketch}[theorem:gauss-sums]
  In classical terms, the product of the basic functions outside $S$ is, up to a constant, the normalized Gauss sum $\tilde{g}_\psi^S(b, c)$: a weighted sum, over the $\tau$ with $\tau^n \mid c$, of the Gauss sums of order $n$ with modulus $c / \tau^n$ and frequency $b$.  The sum over $\tau$ corresponds to the vectors $\tau^n (c, d)$ in the sum over $F^2 - \{0\}$.  The sums are twisted by a character $\chi$ built from characters $\chi_{\mathfrak{p}}$ of $\mathfrak{o}_{\mathfrak{p}}^\times$ at the primes of a finite set $T$, with $\chi_{\mathfrak{p}}^n$ ramified.  These are the simplest twists: at $\mathfrak{p} \in T$ the transform of $1_{\mathfrak{o}_{\mathfrak{p}}^\times} \chi_{\mathfrak{p}}$ is a product of two Gauss sums supported on a single valuation, and the main term vanishes when $T$ is nonempty.  On the dual side the twist becomes $\overline{\chi(c)} \qr{c}{f}_n \qr{f}{c}_n$, and the constant $\varepsilon$, of modulus one, collects the local Gauss sums.  For $n = 1$ the formula reads $\sum_{c \mid b} W_S(c) \chi(c) = \sum_{c \mid b} W_S(b / c) \chi(b / c)$.  Other twists are covered by enlarging $S$ or by changing the frequency $b$ (Remarks \ref{remark:further-conditions} and \ref{remark:frequency-twist}).
\end{sketch}

\begin{sketch}[corollary:quadratic-dual]
  For odd $n$, take $\chi_{\mathfrak{p}} := \eta \kappa^2$ at each prime dividing $h$, where $\eta := \{h, \cdot\}_2$ and $\kappa := \{h, \cdot\}_n$.  Two Hilbert symbols turn the twist into $\chi^{-1} \kappa^2 = \eta$ on the dual side.  One comes from the transform at $\mathfrak{p}$; the other comes from the basic functions at the primes not dividing $h$, and is carried to $\mathfrak{p}$ by the reciprocity law.  On the plane, their product $\{a, u\}_n^2$ is the commutator of two lifted dilations (Remark \ref{remark:lifted-dilations}).  The reciprocity law also leaves the Hilbert symbol $\{h, c\}_S$ at the places in $S$.  For $n = 3$, the twist $\qr{c}{h}_2 \qr{c}{h}_3^2$ is the sextic residue symbol, and the dual twist is quadratic: the exchange seen in the work on the quasi-Riemann hypothesis.
\end{sketch}

% Appendix A

\begin{sketch}[sec:marginals]
  For $n = 1$, the Fourier transform on the plane exchanges integration along the lines $x = \mathrm{const}$ with restriction to the line $x = 0$.  For general $n$, $\mathcal{F}$ exchanges the marginals $m_\Phi(x) = \int_F \Phi_X(x, y) \, d y$ and the axis values $r_\Phi(y) = \Phi_Y(0, y)$, each time by a Fourier transform in an $n$th root variable.  One direction is immediate from the definition of $\mathcal{F}$; the other is the main local computation, and reduces, by the Steinberg relation, to Fourier inversion on $F$.  Together they give the inversion formula, and they are the local input to the proof of the Poisson formula.
\end{sketch}

% Appendix B

\begin{sketch}[sec:strategy]
  The usual proof of Poisson summation applies Fourier inversion on a compact quotient group, but genuine functions are not functions on a group.  Instead, one first treats the pure tensors with a factor of leading term zero and a factor whose transform has leading term zero.  The difference of the two sides, evaluated on the translates of $\Phi$, is invariant under $\mathrm{SL}_2(F)$, and is shown to vanish in three steps:
  \begin{enumerate}
  \item Its constant term vanishes.  Poisson summation on $\mathbb{A}$, applied to marginals and axis values, produces four terms at the origin: two given by the leading terms, which vanish here, and two that cancel by Poisson summation on $\mathbb{A}^\times / \mathbb{A}^{\times n}$.
  \item Hence the difference decays rapidly at the cusp.
  \item $F^2 - \{0\}$ is a single orbit of $\mathrm{SL}_2(F)$ with unipotent stabilizers, so unfolding the integral of the square of the difference and using the first step again gives zero.
  \end{enumerate}
\end{sketch}

\begin{sketch}[sec:boundary]
  For a general $\Phi$, modify the factor at one finite place $w$, without changing its leading term, so that its transform has leading term zero.  Dilation by $\varpi^n$ multiplies $b \Phi_w$ by $q^{n-1}$ and $b \mathcal{F} \Phi_w$ by $q^{n+1}$, so the combination
  \[
    Q_w \Phi_w = \frac{q^{n+1} \Phi_w - \Phi_w(\varpi^n \, \cdot)}{q^{n+1} - q^{n-1}}
  \]
  has $b Q_w \Phi_w = b \Phi_w$ and $b \mathcal{F} Q_w \Phi_w = 0$.  The vanishing case shows that $\Phi \mapsto D(Q \Phi)$ factors through $b$, where $D$ is the difference of the two sides, and $B(b \Phi) := -D(Q \Phi)$ is the boundary functional.
\end{sketch}
